\chapter{Dentro da Máquina: Transformer sem Fumo Sagrado}
\label{ch:transformer}

«Depositei dinheiro no banco» e «sentei-me no banco» pedem interpretações diferentes da mesma palavra. As outras palavras dão contexto. Queremos abrir uma parte do mecanismo que permite combinar essa informação. Não vamos encontrar uma ficha de dicionário com uma pequena pessoa a escolher o significado. Vamos encontrar números, transformações e dependências entre posições.

A primeira leitura pode seguir três passos. Os fragmentos de texto tornam-se listas de números. As camadas combinam informação dessas listas. A saída atribui pontuações às unidades que podem aparecer a seguir. Vamos fazer uma destas contas à mão. A notação geral fica no aprofundamento.

\section{Da lista de números à regra de transformação}

Imagina a lista $(2,3)$ numa ficha de inventário. O primeiro número conta meias e o segundo conta toalhas. Chamamos \emph{vetor} a uma lista ordenada de números deste género. Esta lista tem duas posições, ou duas \emph{dimensões}. Se trocares os números de lugar, mudas a ficha.

Num Transformer de texto, cada identificador de token permite consultar uma linha de uma tabela aprendida. A lista devolvida chama-se \emph{embedding} inicial. Com a tabela fixa, o mesmo identificador recupera a mesma lista. Na ficha de inventário, demos um nome a cada posição. Numa representação aprendida, em geral é a combinação dos números que interessa. O contexto entra nas transformações das camadas seguintes~\citep{vaswani2017}.

Uma \emph{matriz} é uma tabela de números com linhas e colunas. Cada linha pode escrever uma regra. Se a primeira linha é $(1,2)$, multiplica o primeiro número da ficha por 1 e o segundo por 2. Soma os resultados: $1\times2+2\times3=8$. Se a segunda linha é $(0,1)$, obténs $0\times2+1\times3=3$. Aplicar a matriz à ficha $(2,3)$ dá a nova lista $(8,3)$. Cada linha produziu um número da saída.

Este tipo de multiplicação permite que o modelo construa listas diferentes a partir da mesma entrada. O Transformer também usa operações que não se reduzem a multiplicar por uma matriz. Mais adiante, vamos separar o trabalho de cada uma.

\section{Atenção com desenhos de guardanapo}

A \emph{autoatenção}, ou \emph{self-attention}, combina informação de posições da mesma sequência. Pensa nas palavras à volta de «banco». Para formar uma representação daquela posição, o modelo calcula quanto pesa a informação de cada posição disponível. Num modelo que gera da esquerda para a direita, a \emph{máscara causal} indica quais as posições permitidas. Durante o treino, impede o acesso aos tokens seguintes.

De forma simplificada, para cada token temos três vetores:

\begin{itemize}[leftmargin=*]
\item q — \emph{query}: “o que eu procuro”.
\item k — \emph{key}: “o que eu ofereço”.
\item v — \emph{value}: “a informação que posso fornecer”.
\end{itemize}

As frases entre aspas são mnemónicas. Queries, keys e values são projeções calculadas a partir das representações, não intenções de pequenos agentes. «Procuro» não significa uma pesquisa num catálogo externo nem atribui vontade ao token. A analogia ajuda a ler o cálculo e termina aí.

Comparamos uma query com cada key. Um \emph{produto interno} multiplica os números que estão na mesma posição e soma os produtos, como fizemos com uma linha da matriz. O resultado, depois de um ajuste de escala, é uma pontuação. Para combinar os valores, precisamos de transformar as pontuações em pesos que somem um. A regra que faz essa transformação chama-se \emph{softmax}.

Supõe que os pesos de duas posições são um quarto e três quartos, e que os dois valores são 2 e 6. Dás a cada valor a sua parte: $\tfrac14\times2+\tfrac34\times6=5$. É a média ponderada da página de \hyperref[sec:ler-numeros]{números e fórmulas}. A mesma conta pode ser aplicada a cada componente se os valores forem listas. Os números são didáticos.

\textbf{Aprofundamento opcional.} $Q$, $K$ e $V$ agrupam, por linha, as queries, keys e values das posições. A fórmula escreve a comparação, a conversão em pesos e a combinação:

\[
\operatorname{Attention}(Q,K,V)=
\operatorname{softmax}\!\left(\frac{QK^{\mathsf T}}{\sqrt{d_k}}+M\right)V.
\]
Com $n$ posições, $Q$ e $K$ têm $n$ linhas e $d_k$ colunas. $V$ tem $n$ linhas e $d_v$ colunas. O $\mathsf T$ em cima de $K$ indica a \emph{transposição}, que troca as suas linhas e colunas. Cada entrada de $QK^{\mathsf T}$ compara uma query com uma key. O símbolo $\sqrt{d_k}$ representa a raiz quadrada, como 2 para $d_k=4$. Dividir por esse valor controla a escala. A softmax é aplicada a cada linha, uma por query~\citep{vaswani2017}.

$M_{ij}$ é a entrada da máscara na linha $i$, coluna $j$. Para a máscara causal, vale zero quando $j\leq i$ e $-\infty$ quando $j>i$. Isto permite as posições anteriores e a atual. O símbolo $-\infty$ representa um limite que dá peso zero às futuras. O capítulo usa uma implementação didática com apenas as posições permitidas.

Numa linha com pontuações permitidas $[0,\ln 3]$, a softmax calcula $e^0=1$ e $e^{\ln 3}=3$. Aqui, $e$ é uma constante próxima de $2{,}718$ e $\ln 3$ é o expoente a que a elevamos para obter 3. Divide por $1+3$ e obténs os pesos $[1/4,3/4]$ usados acima. Somar a mesma constante às duas pontuações multiplica os dois valores exponenciados pelo mesmo fator. A divisão pela soma cancela esse fator. O programa de fundamentos subtrai o máximo antes de exponenciar, o que evita elevar $e$ a um número positivo muito grande.

\section{Uma máscara impede uma pista vinda do futuro}
Imagina agora três posições. Estamos a calcular a segunda e os dois valores acessíveis são os do exemplo, dois e seis. A terceira posição contém, na construção didática, o valor mil. Num descodificador causal, essa posição futura não pode participar. A máscara atribui-lhe peso zero, em vez de a deixar alterar a média.

Na fórmula, o score dessa posição recebe $-\infty$ antes da softmax, tornando nula a sua contribuição na interpretação matemática. Uma implementação concreta precisa de uma representação e de operações apropriadas. Para repetir a nossa conta escalar, basta usar apenas os dois scores permitidos e os dois valores correspondentes; o helper didático não é uma biblioteca de máscaras tensoriais.

O cuidado tem uma razão de treino. Se uma posição pudesse consultar o token futuro que se pretende prever, o sistema receberia uma pista que não existirá na geração. É outra forma de reconhecer a fuga de informação que investigámos no capítulo~\ref{ch:ml}. O treino em várias posições pode decorrer em paralelo sem lhes conceder acesso ao futuro.

O que isto faz, em termos de desenho de guardanapo:

\begin{enumerate}[leftmargin=*]
\item Cada posição compara a sua query (Q) com keys (K) das posições permitidas.
\item Calcula “graus de afinidade” (os scores) com cada um.
\item Passa esses scores por uma softmax para obter pesos que somam 1.
\item Faz uma média ponderada das \emph{values} V usando esses pesos.
\end{enumerate}

No fim, a representação que resulta da atenção combina os valores das posições disponíveis. Um peso maior dá a uma posição mais influência nessa combinação.

Do ponto de vista de prompts:

\begin{itemize}[leftmargin=*]
\item A posição de informação relevante pode alterar o desempenho observado.
\item Instruções contraditórias tornam ambíguo o comportamento esperado; uma resposta não revela por si só a distribuição de atenção interna.
\item Exemplos estruturados ajudam a especificar a tarefa, mas o seu efeito deve ser medido.
\end{itemize}

Mais tarde, quando falarmos de \emph{janela de contexto} e de “lost in the middle”, vais ver como isto se traduz em consequências práticas para o design de mensagens longas.

\section{Várias cabeças e várias camadas não são vários especialistas humanos}
Uma cabeça de atenção usa as suas projeções para produzir uma combinação. A atenção multicabeça permite várias dessas operações, cujos resultados são concatenados e novamente projetados. A arquitetura oferece canais de cálculo diferentes; não lhes atribui por definição profissões como «gramático» e «detetor de factos».

Os parâmetros aprendidos podem levar a padrões interessantes, que exigem investigação para serem interpretados. Dar um nome a uma cabeça depois de observar alguns exemplos não demonstra uma função exclusiva em todos os contextos. Também não permite deduzir a verdade da saída a partir de um peso de atenção isolado.

Camadas sucessivas recebem representações já transformadas. Ligações residuais preservam caminhos de informação e as transformações feed-forward acrescentam operações sobre as representações. A saída final depende do conjunto. Por isso o desenho de uma única cabeça é um modelo de leitura do mecanismo, não uma explicação completa de uma resposta específica.

\begin{promptlabbox}{Ver o impacto da ordem das instruções}
\textbf{Objetivo e contexto:} Comparar organização de instruções com o mesmo parágrafo, sem mudar os requisitos.

Experimenta comparar estes dois prompts:

\medskip

\textbf{Prompt A}

\begin{lstlisting}
Quero que analises o seguinte texto.

[COLOCA AQUI UM PARÁGRAFO QUALQUER]

No fim, faz:

Um resumo em 3 pontos.

Uma crítica em tom cínico.

Uma sugestão de melhoria.

Começa pelo resumo.
\end{lstlisting}

\medskip

\textbf{Prompt B}

\begin{lstlisting}
Quero que faças três coisas, por esta ordem:

Um resumo em 3 pontos.

Uma crítica em tom cínico.

Uma sugestão de melhoria.

Só depois disso é que vais ler o seguinte texto
e aplicar os passos anteriores:

[COLOCA AQUI O MESMO PARÁGRAFO]

Começa pelo resumo.
\end{lstlisting}

Compara:
\begin{itemize}[leftmargin=*]
\item Em qual deles o modelo segue melhor a ordem das tarefas?
\item Em qual deles parece “esquecer-se” de algum passo?
\end{itemize}

Este teste compara comportamento observável. Não revela pesos de atenção nem prova uma explicação causal sobre os mecanismos internos do modelo.

\textbf{Resultado esperado e validação:} Assinala se há três pontos de resumo, uma crítica e uma sugestão, pela ordem pedida. Repete em sessões separadas; empate é um resultado admissível.

\textbf{Falhas possíveis:} Mudar simultaneamente texto, tarefa e ordem; interpretar o resultado como medição direta de atenção.
\end{promptlabbox}

\section{Embeddings e espaço latente}

Quando dizemos que o modelo “entende” que \texttt{gato} e \texttt{cão} são mais parecidos entre si do que com \texttt{parafuso}, não estamos a atribuir-lhe intuição animal. Estamos a falar de \emph{embeddings}: representações vetoriais aprendidas.

Durante o treino, o objetivo e os dados dão origem a regularidades nas representações. A geometria depende do modelo, da camada e da medida de proximidade. Num \emph{espaço latente}, podemos encontrar padrões como:

\begin{itemize}[leftmargin=*]
\item Palavras com usos semelhantes podem ter representações próximas; antónimos também podem partilhar contextos, pelo que proximidade não significa sinonímia.
\item Palavras que partilham função gramatical ou papel semântico têm padrões reconhecíveis.
\item Certos embeddings estáticos apresentam regularidades relacionais aproximadas, estudadas, por exemplo, por Mikolov e colaboradores~\citep{mikolov2013}. Analogias vetoriais não são identidades algébricas nem uma regra geral para representações contextuais.
\end{itemize}

Porque é que isto interessa para prompts?

\begin{itemize}[leftmargin=*]
\item Mudar o vocabulário altera a sequência de tokens e pode mudar a resposta; o efeito não se deduz só de uma distância entre embeddings.
\item Exemplos de estilo acrescentam padrões ao contexto. Compara a sua utilidade em respostas novas, sem lhes atribuir uma direção geométrica que não mediste.
\item Instruções positivas e negativas são alternativas a testar. «Escreve como se estivesses a falar com um colega» torna um objetivo observável; não garante que seja superior em todas as tarefas.
\end{itemize}

\begin{keyideabox}{Prompts empurram o contexto no espaço latente}
Cada palavra que usas no prompt empurra o contexto para uma região diferente do espaço latente. É uma metáfora útil, mas não existe um único «embedding do contexto» que possas deslocar como uma peça num tabuleiro. Formulações distintas condicionam o cálculo; a sua eficácia deve ser medida na tarefa.
\end{keyideabox}

\section{O que existe para além da atenção}
A entrada começa por identificadores de tokens e embeddings. Informação de posição permite distinguir «o cão mordeu o homem» de «o homem mordeu o cão». Depois de combinar informação de várias posições, o bloco aplica também transformações \emph{feed-forward}. Estas operam na lista de cada posição. Uma \emph{ligação residual} soma o resultado de uma operação à entrada dessa operação, criando um caminho direto para a informação anterior. A \emph{normalização} ajusta a escala dos números em etapas do cálculo. A composição e a ordem têm variantes. Na saída, uma transformação calcula pontuações sobre o vocabulário.

No treino, conhecemos a sequência de referência e calculamos perdas em várias posições em paralelo, respeitando a máscara. Na geração autoregressiva, acrescentamos tokens sucessivamente. Uma cache de keys e values pode evitar recalcular parte do histórico; não é uma memória permanente de conversas.

\begin{figure}[htbp]
\centering
\begin{tikzpicture}[node distance=7mm, every node/.style={draw, rounded corners=1pt, align=center, text width=35mm, minimum height=9mm, font=\small}, >=Stealth]
\node (in) {Tokens e posições};
\node (att) [below=of in] {Atenção causal};
\node (ff) [below=of att] {Feed-forward, normalização e resíduos};
\node (out) [below=of ff] {Logits e descodificação};
\draw[->] (in)--(att); \draw[->] (att)--(ff); \draw[->] (ff)--(out);
\end{tikzpicture}
\caption{Percurso simplificado num Transformer de texto com descodificador causal. Os blocos repetem-se; não é uma arquitetura comercial específica.}
\label{fig:transformer}
\end{figure}

Uma alteração no prompt muda a entrada deste cálculo. Não te dá acesso direto aos parâmetros, às posições internas ou a um interruptor de «verdade». O desenho é útil precisamente porque mantém essas fronteiras visíveis.
\section{Como os parâmetros aprendem}
Uma rede começa com parâmetros e um objetivo de treino. Imagina que nos dados a palavra seguinte é «meia». Um modelo atribui-lhe 90\% de probabilidade e outro atribui-lhe 10\%. A regra de treino penaliza mais a segunda previsão. Repete a comparação em muitas posições e ajusta os parâmetros na direção que reduz a perda, como na reta do capítulo~\ref{ch:ml}.

\textbf{Aprofundamento opcional.} Uma perda comum é a média da log-probabilidade negativa dos tokens de referência:
\[
\mathcal{L}(\theta)=-\frac{1}{N}\sum_{k=1}^{N}
\log P_{\theta}(t_k\mid t_{<k}).
\]
Aqui, $\theta$ (teta) representa os parâmetros, $N$ conta as posições avaliadas e $t_{<k}$ reúne os tokens anteriores à posição $k$. O sinal negativo faz com que a perda cresça quando a probabilidade do token correto baixa. Com logaritmo natural, $-\ln(0{,}9)$ é cerca de $0{,}11$ e $-\ln(0{,}1)$ é cerca de $2{,}30$. Esta conta mede o ajuste aos exemplos.

O \emph{gradiente} reúne as inclinações da perda quando mexemos em cada parâmetro. O símbolo $\nabla_{\theta}\mathcal{L}$ representa essa lista. Uma atualização simples escreve-se $\theta\leftarrow\theta-\eta\nabla_{\theta}\mathcal{L}$. A seta indica que substituímos os valores antigos pelos novos e $\eta$ (eta) define o tamanho do passo. Os otimizadores, programas que escolhem os passos de ajuste, podem usar variantes. A retropropagação permite calcular essas inclinações através das camadas~\citep{goodfellow2016}.

Treinar exige muitos exemplos, computação e avaliação. Durante a inferência normal, não repetimos esta atualização por cada mensagem: usamos parâmetros fixos para calcular uma distribuição condicionada ao contexto. É por isso que fornecer uma política no prompt e alterar os pesos são operações diferentes.

\section{Janela de contexto e “lost in the middle”}

Um LLM não lê texto infinito. Ele trabalha com uma \emph{janela de contexto} finita: um número máximo de tokens que consegue considerar de uma vez. Os limites variam por modelo e serviço; confirma separadamente capacidade de entrada, limite de geração e tratamento do histórico. Uma interface pode rejeitar uma entrada longa, resumi-la ou truncá-la; não há uma regra universal.

Dois problemas práticos aparecem aqui:

\begin{enumerate}[leftmargin=*]
\item \textbf{Material omitido ou pedido rejeitado}: ao ultrapassar um limite, a aplicação pode rejeitar a entrada ou cortar parte dela. O modelo não utiliza diretamente o trecho que não recebeu; a estratégia de corte não é universal.
\item \textbf{Lost in the middle}: Liu e colaboradores observaram degradação associada à posição da evidência nas tarefas e modelos que testaram~\citep{liu2024}. É um resultado experimental importante, não uma lei invariável de todos os modelos com contexto longo.
\end{enumerate}

Para prompting, isto significa:

\begin{itemize}[leftmargin=*]
\item Uma organização a testar coloca instruções na camada apropriada e compara posições da evidência, conservando os requisitos.
\item Secções e marcadores podem facilitar seleção e verificação de material; o efeito sobre o modelo deve ser avaliado, sem os tratar como garantia de acerto.
\item Muitas conversas longas com “pequenos ajustes” acabam num estado em que ninguém (nem tu, nem o modelo) sabe exatamente que instruções ainda estão ativas.
\end{itemize}

Mais à frente, quando falarmos de \emph{prompt chaining} e de RAG, vamos ver como contornar isto com pipelines em vez de tentar meter o mundo inteiro numa única janela.

\section{RLHF, RLAIF e camadas de alinhamento}

Até agora falámos do modelo base enquanto preditor de texto. Mas o que encontras numa interface pública (tipo chat) não é só o modelo base; é um modelo que pode ter recebido pós-treino para seguir instruções e políticas. Alinhamento tem aceções mais amplas do que cortesia: precisamos de explicitar os comportamentos e objetivos avaliados.

Uma receita histórica importante, a do RLHF com modelo de recompensa, é~\citep{ouyang2022}:

\begin{enumerate}[leftmargin=*]
\item Treinar um modelo base em montes de texto para prever o próximo token.
\item Realizar ajuste supervisionado com demonstrações de respostas; recolher depois comparações humanas entre respostas para treinar preferências.
\item Treinar um \emph{modelo de recompensa} que estime preferências a partir das comparações recolhidas.
\item Ajustar o modelo após o treino supervisionado com um método de reforço orientado pela recompensa, na receita histórica descrita.
\end{enumerate}

Isto não é uma receita única. O ajuste supervisionado usa exemplos de respostas; RLAIF usa avaliações de IA~\citep{bai2022}; DPO otimiza preferências sem exigir a mesma sequência explícita de modelo de recompensa e otimização por reforço~\citep{rafailov2023}. Os fornecedores podem combinar métodos.

O resultado pretendido é um modelo mais:

\begin{itemize}[leftmargin=*]
\item educado,
\item menos tóxico,
\item útil dentro de certos limites,
\item capaz de seguir as políticas avaliadas.
\end{itemize}

Do ponto de vista de prompting, isto explica várias coisas que vês no dia a dia:

\begin{itemize}[leftmargin=*]
\item Respostas cheias de disclaimers e avisos legais.
\item Recusa em entrar em certos temas, mesmo quando a pergunta é académica.
\item Tendência para dizer “como modelo de IA, não posso…” em situações em que, tecnicamente, até podia responder.
\end{itemize}

Quando mais tarde falarmos de jailbreaks e DAN, estás na verdade a tentar contornar estas camadas de alinhamento. Quando falarmos de segurança e ética, vamos ver como usar esse conhecimento para \emph{defender} sistemas em produção, em vez de os rebentar.

\begin{warningboxenv}{Não confundir alinhamento com competência}
Um modelo pode ser altamente “alinhado” segundo determinados objetivos e ao mesmo tempo alucinar factos básicos. Alinhamento não é sinónimo de precisão. Quando desenhas prompts, tens de lidar com estas duas dimensões em separado.
\end{warningboxenv}

O motor já tem peças reconhecíveis. Falta perceber a unidade que atravessa todo o percurso: o token. O seu tamanho e a maneira como o texto é dividido afetam limites e custos, mesmo quando a frase parece curta a quem a lê.

